\section*{Hoofdstuk \ref{ch:b3:continuum-elasticity} --- Continuümmechanica en elasticiteit}

\begin{solution}{exo:b3:continuum-elasticity:1}
(a) $\unit{Pa} = \unit{N/m^2} = \unit{kg.m^{-1}.s^{-2}}$; rek is een
zuiver getal. (b) $\varepsilon = \sigma/E = \num{e-3}$. (c) De lengte
verdubbelen is $\varepsilon = 1$: duizend keer voorbij ``klein'', en de
reactie van rubber is daar sterk niet-lineair (en van een andere,
entropische oorsprong) --- de \omterm{thm:b3:continuum-elasticity:hooke}{wet van Hooke} opent de kromme slechts. (d)
$w =
\sigma^2/2E$: staal $(\num{2e8})^2/2(\num{2e11}) =
\qty{1e5}{J/m^3}$; rubber $(\num{e6})^2/2(\num{2e6}) =
\qty{2.5e5}{J/m^3}$ --- het zachte materiaal slaat bij werkspanning
\emph{meer} per eenheid van volume op, en daarom zijn katapulten van
rubber en niet van staal.
\end{solution}

\begin{solution}{exo:b3:continuum-elasticity:2}
(a) $\varepsilon_{ij} = \alpha\delta_{ij}$: isotrope uitzetting,
$\delta V/V = 3\alpha$. (b) $\varepsilon_{xy} = \varepsilon_{yx} =
\gamma/2$, de rest nul: zuivere afschuiving, volume onveranderd. (c)
$\varepsilon_{ij} = 0$: starre rotatie, geen vervorming. (d)
$\operatorname{diag}(\varepsilon, -\nu\varepsilon, -\nu\varepsilon)$:
de trekproef --- rek langs $x$, poissoncontractie dwars daarop.
\end{solution}

\begin{solution}{exo:b3:continuum-elasticity:3}
(a) Met alleen $\sigma_{zz}(z)$ en het gewicht: $\partial_z\sigma_{zz} =
\rho g$ (met druk positief naar beneden), $\sigma_{zz} = \rho
gh$ aan de
voet. (b) $h_{\max} = \sigma_{\text{c}}/\rho g =
\num{2e8}/(2700 \times 9.81) \approx \qty{7.6}{km}$. (c) De Mount
Everest zit op de verbrijzelingsgrens van zijn eigen voet --- de bergen
van de aarde zijn zo hoog als de sterkte van gesteente toelaat, en niet
hoger. (d) Een kegel met hoogte $h$ belast zijn grondvlak met slechts
$\rho gh/3$ (een derde van de zuil: de massa groeit mee met de
doorsnede), zodat een bergvormige stapel ongeveer drie keer zo hoog kan
staan als een zuil.
\end{solution}

\begin{solution}{exo:b3:continuum-elasticity:4}
(a) $-\nu\varepsilon = \num{-3e-4}$. (b) $\delta V/V = (1 -
2\nu)\varepsilon = \num{4e-4}$. (c) $\nu = 1/2$: onsamendrukbare
vervorming --- rubber. (d) Voor $\nu > 1/2$ zou de compressiemodulus $K =
E/3(1 - 2\nu)$ negatief zijn: van alle kanten samendrukken zou het volume
doen \emph{groeien}, en het materiaal zou energie vrijmaken door in te
storten --- geen stabiele vaste stof kan dat.
\end{solution}

\begin{solution}{exo:b3:continuum-elasticity:5}
(a) Op een halve cilinder van eenheidslengte duwt de druk met $p
\times 2R$ (geprojecteerd oppervlak), terwijl de twee doorgesneden wanden
met $2\sigma_\theta t$ terugtrekken: $\sigma_\theta = pR/t$. (b) Op een
dwarsdoorsnede geldt $p\pi R^2 = \sigma_z\,2\pi Rt$: $\sigma_z = pR/2t$
--- de helft. Een worstje barst \emph{in de lengte} open: de grotere
ringspanning scheurt het vel langs de as. (c) $\sigma_\theta = \num{2e7} \times
0.09/0.005 = \qty{360}{MPa}$: een factor $2$ onder de vloeigrens
--- en daarom worden flessen beproefd en gekeurd. (d) Een bol draagt
$pR/2t$ in elke richting: halfbolvormige einden halveren de spanning en
vermijden de hoeken waar vlakke einden haar zouden opstapelen.
\end{solution}

\begin{solution}{exo:b3:continuum-elasticity:6}
(a) $G = E/2(1 + \nu) = \qty{1.0}{MPa}$. (b) $\sigma_{xy} = F/S =
150/\num{e-2} = \qty{15}{kPa}$; $\gamma = \sigma_{xy}/G = 0.015$;
verplaatsing $\gamma h = \qty{0.3}{mm}$. (c) Een rubbernetwerk verandert
van \emph{vorm} door zijn ketens louter anders te richten (makkelijk),
maar van \emph{volume} veranderen betekent de ketens dichter opeen
pakken, even moeilijk als een vloeistof samendrukken: $G \ll K$, dus $\nu \to 1/2$. (d) De zachtheid bij afschuiving: de oplegging laat de bodem
er horizontaal onder schudden en geeft daarbij weinig kracht door, maar
blijft onder druk stijf genoeg om het gewicht van het gebouw te dragen.
\end{solution}

\begin{solution}{exo:b3:continuum-elasticity:7}
(a) $k = ES/L = \num{1.2e9} \times \num{8e-5}/20 = \qty{4.8}{kN/m}$.
(b) $mg(h + x) = \tfrac12 kx^2$: $2400x^2 - 785x - 3140 = 0$, $x =
\qty{1.3}{m}$. (c) $F = kx \approx \qty{6.3}{kN}$; vertraging $F/m
\approx \qty{79}{m/s^2} \approx 8g$ --- de reden dat touwen met opzet
rekbaar worden gemaakt. (d) Een deel van de opgenomen energie ging naar
het breken van vezels en plastische herschikking: het touw zit nu op een
aangetaste spanning--rekkromme, stijver en zwakker, en de volgende val
zou in beide opzichten harder aankomen.
\end{solution}

\begin{solution}{exo:b3:continuum-elasticity:8}
(a) De top van een buisje met straal $r$ draait over de boog $r\theta$
over de lengte $L$: $\gamma = r\theta/L$. (b) $\Gamma = \int_0^a r\,(G r\theta/L)\,
2\pi r\,\dd r = (\pi Ga^4/2L)\,\theta$. (c) Met $2^4 = 16$. (d) $C =
\pi \times \num{4e10} \times (\num{2.5e-5})^4/2 =
\qty{2.5e-8}{N.m/rad}$; $T = 2\pi\sqrt{I/C} = 2\pi\sqrt{\num{5e-4}/
\num{2.5e-8}} \approx \qty{900}{s}$ --- een kwartier per slingering:
gevoelig genoeg om de zwaartekracht van loden bollen te voelen.
\end{solution}

\begin{solution}{exo:b3:continuum-elasticity:9}
(a) $\mu = E/2(1+\nu) = \qty{78}{GPa}$, $\lambda = E\nu/(1+\nu)(1-2\nu)
= \qty{107}{GPa}$: $c_{\text{P}} = \sqrt{262\times10^9/7850} =
\qty{5.8}{km/s}$, $c_{\text{S}} = \qty{3.1}{km/s}$. (b) De dunne staaf
geeft $\sqrt{E/\rho} = \qty{5.0}{km/s}$: in de staaf mogen de zijkanten
vrij uitpuilen (poissonontspanning), wat de reactie zachter maakt; in het
volle materiaal verbiedt de omringende materie dat. (c) $c_{\text{S}} = 0$; $c_{\text{P}} = \sqrt{\num{2.2e9}/1000} = \qty{1.5}{km/s}$ --- de
geluidssnelheid van water. (d) P beweegt de bodem langs de straal (bijna
verticaal bij een diepe bron), S dwars daarop: de drie componenten van
een seismometer scheiden de polarisaties, en de P/S-vertraging dateert
vervolgens de afstand.
\end{solution}

\begin{solution}{exo:b3:continuum-elasticity:10}
(a) Een boog op straal $R + y$ heeft lengte $(R + y)\alpha$ tegen
$R\alpha$ op het neutrale vlak: $\varepsilon = y/R$, $\sigma =
Ey/R$. (b) $M = \int y\sigma\,\dd S = (E/R)\int y^2\,\dd S = EI/R$;
voor de rechthoek is $I = bh^3/12$. (c) Plat: $I = bh^3/12$; op zijn
kant: $hb^3/12$ --- verhouding $(b/h)^2 = 25$ voor $b/h = 5$: hetzelfde
hout, vijfentwintig keer stijver, louter door de meetkunde. (d) $I = 0.50 \times
0.035^3/12 = \qty{1.8e-6}{m^4}$, $EI = \qty{2.1e4}{N.m^2}$;
$\delta =
FL^3/3EI = 736 \times 1.8^3/(3 \times \num{2.1e4}) \approx
\qty{7}{cm}$: precies de prettige veerkracht van een duikplank.
\end{solution}

\begin{solution}{exo:b3:continuum-elasticity:11}
(a) In de uitgeweken stand werkt de axiale last $P$ met arm $y(x)$ om de
doorsnede in $x$: $M = -Py$, en $M = EIy''$ geeft $EIy'' = -Py$. (b) $y = A\sin(x\sqrt{P/EI})$ met $y(L) = 0$: een $A$ ongelijk nul is pas mogelijk
bij $\sqrt{P/EI}\,L = \pi$, dus $P_{\text{c}} =
\pi^2EI/L^2$; daaronder
is de rechte zuil de enige oplossing, daarboven kost uitwijken geen
kracht. (c) $I = 0.03 \times 0.0015^3/12 =
\qty{8.4e-12}{m^4}$:
$P_{\text{c}} = \pi^2 \times 12\times10^9 \times
\num{8.4e-12}/1^2 \approx \qty{1}{N}$ --- het gewicht van een appel: een meetlat knikt
onder één vinger. (d) $I$ groeit naarmate materiaal verder van de as
ligt: een buis legt haar hele doorsnede bij grote $y$ en maximaliseert
zo $I$ (en dus $P_{\text{c}}$) bij gegeven materiaalgewicht --- het
ontwerp van botten, bamboe en fietsframes.
\end{solution}

\begin{solution}{exo:b3:continuum-elasticity:12}
(a) Eén keten per oppervlak $a^2$; rekken met rek $\varepsilon$ rekt elke
veer met $\delta = \varepsilon a$, geeft een kracht $\kappa
\varepsilon a$
per keten en een spanning $\kappa\varepsilon/a$: $E =
\kappa/a$. (b)
$\rho = m/a^3$, dus $E/\rho = \kappa a^2/m$ en $c =
a\sqrt{\kappa/m}$.
(c) $c = \num{2.5e-10}\sqrt{50/\num{e-25}} \approx
\qty{5.6}{km/s}$. (d)
Precies de gemeten paar km/s van metalen: de stijfheid van alles wat vast
is --- en de snelheid van elke aardbevingsgolf --- is de stijfheid van de
chemische binding.
\end{solution}

\begin{solution}{pb:b3:continuum-elasticity:1}
\textbf{1.} $\mu = E/2(1+\nu) = \qty{30}{GPa}$; $\lambda =
E\nu/(1+\nu)(1-2\nu) = \qty{30}{GPa}$: voor $\nu = 1/4$ is $\lambda =
\mu$.
\textbf{2.} $c_{\text{P}} = \sqrt{3\mu/\rho} = \sqrt{\num{9e10}/2700}
= \qty{5.8}{km/s}$; $c_{\text{S}} = \sqrt{\mu/\rho} =
\qty{3.3}{km/s}$; verhouding $\sqrt3$.
\textbf{3.} $\varepsilon \sim 2/\num{5e4} = \num{4e-5}$; $\sigma =
E\varepsilon \approx \qty{3}{MPa}$ --- de gemeten ``spanningsval'' van
echte aardbevingen bedraagt inderdaad een paar MPa.
\textbf{4.} $w = \tfrac12 E\varepsilon^2 \approx \qty{60}{J/m^3}$;
over $\num{5e13}{}\,\unit{m^3}$: $\sim\qty{3e15}{J}$, ongeveer
driekwart megaton --- een zware aardbeving.
\textbf{5.} Van onderen beweegt P de bodem verticaal --- gebouwen zijn
gebouwd om verticale lasten te dragen; S beweegt hem \emph{horizontaal},
de richting waarin gebouwen het zwakst zijn.
\textbf{6.} Alleen de \omterm{prop:b3:continuum-elasticity:waves}{P-golf}: dat is een drukgolf; de \omterm{prop:b3:continuum-elasticity:waves}{S-golf} is zuivere
afschuiving, bij constant volume.
\textbf{7.} $t_{\text{P}} = d/c_{\text{P}}$, $t_{\text{S}} =
d/c_{\text{S}}$: $\Delta t = d(1/c_{\text{S}} - 1/c_{\text{P}})$,
omgekeerd zoals aangegeven.
\textbf{8.} $1/(1/3.33 - 1/5.77)\,\unit{km/s} \approx
\qty{7.9}{km}$ per seconde vertraging --- de ``acht kilometer per
seconde'' van de seismoloog in het veld.
\textbf{9.} $d \approx 7.9 \times 28 \approx \qty{220}{km}$.
\textbf{10.} Drie (in het algemeen): elk station kent een cirkel van
mogelijke epicentra; twee cirkels snijden elkaar in twee punten, en het
derde beslist.
\textbf{11.} $t_{\text{P}} = 80/5.8 = \qty{14}{s}$, $t_{\text{S}} =
80/3.3 = \qty{24}{s}$: ongeveer \qty{10}{s} waarschuwing tussen de schok
en het verwoestende schudden.
\textbf{12.} S komt aan op $200/3.3 = \qty{60}{s}$. De P bereikt
\qty{20}{km} na \qty{3.5}{s}; een radiobericht is praktisch ogenblikkelijk,
dus kan de stad bijna een minuut waarschuwing krijgen --- huidige
systemen halen tientallen seconden.
\textbf{13.} $2R_{\text{E}}/c \approx 12742/10 \approx \qty{1270}{s}$,
zo'n twintig minuten.
\textbf{14.} Dat \omterm{prop:b3:continuum-elasticity:waves}{elastische golven} haar überhaupt doorkruisen: de diepe
aarde is niet door en door gesmolten --- zij geleidt, en (bij S-golven
door de mantel) schuift zelfs af als een vaste stof.
\textbf{15.} Twee stralen en de hoek $\Delta$ ertussen: de koorde is
$2R_{\text{E}}\sin(\Delta/2)$.
\textbf{16.} $2 \times 6371 \times \sin30^\circ = \qty{6371}{km}$;
$t \approx \qty{640}{s} \approx \qty{11}{min}$.
\textbf{17.} Voorbij $103^\circ$ moet elke straal door een diep centraal
gebied; is dat gebied \emph{vloeibaar} ($\mu = 0$), dan draagt het geen
schuifgolf: de S-golven sterven daar. Vloeistoffen zijn de
aggregatietoestand die geen afschuiving kan weerstaan.
\textbf{18.} Een vloeistof draagt nog wel samendrukking: P-golven
doorkruisen de kern, maar trager --- zij breken dus scherp aan haar rand,
en de gebogen stralen laten een ringvormige schaduw achter in plaats van
een scherpe snede.
\textbf{19.} In het middelpunt van de aarde zit een vloeibare kern.
\textbf{20.} De kleinste afstand van de koorde tot het middelpunt is
$R_{\text{E}}\cos(\Delta/2)$; scheren betekent
$R_{\text{E}}\cos(\Delta^*/2) = R_{\text{c}}$, dus $\Delta^* =
2\arccos(R_{\text{c}}/R_{\text{E}})$.
\textbf{21.} $R_{\text{c}} = 6371\cos(51.5^\circ) \approx
\qty{3970}{km}$.
\textbf{22.} $+14\%$ te groot. De snelheid groeit met de diepte, dus
buigen echte stralen naar het oppervlak terug: een straal waarvan het
diepste punt de kern net raakt komt onder een \emph{kleinere} hoek boven
dan de rechte koorde door dezelfde diepte --- dus legt de omkering met
rechte stralen het raakpunt vanuit de waargenomen $103^\circ$ te ondiep,
en overschat zij de kern.
\textbf{23.} Zwakke P-energie binnen de schaduw betekent dat iets binnen
de vloeibare kern stralen weer naar buiten buigt: een afzonderlijke
\emph{binnenkern} (Lehmann, 1936).
\textbf{24.} Schuifgolven bestaan alleen in vaste stoffen: een \omterm{prop:b3:continuum-elasticity:waves}{P-golf} die
in de binnenkern in een schuifgolf overgaat en weer terug (de fase die
PKJKP heet) zou haar vastheid bewijzen --- de detectie ervan, lang
gezocht, is het aanvaarde bewijs.
\textbf{25.} Twee snelheden in gesteente ($5.8$ en \qty{3.3}{km/s}), één
ontbrekende golf voorbij $103^\circ$ en de meetkunde van de koorde geven
een vloeibare kern van $\approx\qty{4000}{km}$ --- binnen $15\%$ van de
moderne \qty{3480}{km}: de \omterm{thm:b3:continuum-elasticity:hooke}{wet van Hooke}, op planetaire schaal gelezen.
\end{solution}
